% STATUS: reviewed
\chapter{Terminų žodynas ir santrumpos}\label{app:glossary}

Čia surinkti žodžiai, kuriais šis dokumentas vadina pačios sistemos dalykus, užrašai, kuriuos
konsolė abiem savo kalbomis rodo kiekvienai būsenai, ir teisės aktų bei institucijų santrumpos,
kurių skaitytojas ne iš Lietuvos gali nežinoti.

\section{Sistemos terminai}\label{sec:glossary:terms}

\begin{xltabular}{\textwidth}{@{}L{36mm}Y@{}}
\toprule
\hd{Terminas} & \hd{Reikšmė šioje sistemoje} \\
\midrule
\endfirsthead
\toprule
\hd{Terminas} & \hd{Reikšmė šioje sistemoje} \\
\midrule
\endhead
\bottomrule
\endlastfoot
Apimties peržiūra (\en{scope review}) & Sustojimas ataskaitos rengimo viduryje, kai žmogus
  patvirtina teisės aktų sąrašą, prieš pradedant brangiąją dalį. \\
Ataskaita (\en{report}) & Vienos įmonės atitikties ataskaita, parengta tam tikru metu. \\
DI tinklų sąsaja (\en{gateway}) & Fakulteto DI tinklų sąsaja: aprašytame diegime vienintelis
  kelias, kuriuo sistema pasiekia kalbos modelį. \\
Dokumentų rinkinys (\en{library}) & Visi vienos įmonės dokumentai. \\
Fragmentas (\en{chunk}) & Teisės akto ar įmonės dokumento ištrauka, saugoma kartu su jos
  vektoriumi paieškai. \\
Instrukcija, instrukcijos versija (\en{prompt, prompt version}) & Saugomi nurodymai vienos rūšies
  modelio iškvietimui ir viena nekeičiama jų versija. \\
Inventorius (\en{inventory}) & Tai, ką nustato vienas įmonės dokumentas, vieną kartą perskaitytas
  ir užrašytas struktūrizuotu įrašu: kontrolės priemonės su citatomis, rolės, registrai, patvirtinimas,
  nuostatos, kurias jis cituoja. \\
Išvada (\en{finding}) & Vienos pareigos rezultatas vienoje ataskaitoje: būsena, vertinimas,
  cituojami įrodymai, rekomendacija, pasekmė ir, jei peržiūrėtojas ją įrašė, priežastis. \\
Katalogas (\en{catalogue}) & Teisės aktų, su kuriais produktas pateikiamas, sąrašas ir jų
  šaltiniai. \\
Kelio kortelė (\en{path card}) & Šiame dokumente: dešimties eilučių lentelė, aprašanti vieną
  vietą, kur iškviečiamas modelis. \\
Nuorodų patikra (\en{citation check}) & Dokumente cituojamo straipsnio patikra: ar tokia
  nuostata egzistuoja ir ar ji reglamentuoja tai, ką dokumentas jai priskiria. Ataskaitoje tai
  vadinama \emph{teisės aktų nuorodų patikra}. \\
Nuostata, vienetas (\en{provision, unit}) & Teisės akto straipsnis, dalis, punktas ar sąlyga;
  kaip saugoma eilutė --- vienetas. \\
Pareiga (\en{obligation}) & Teisės akto nustatyta prievolė, vieną kartą užrašyta nurodant, kas ją
  turi vykdyti, kokiai sąlygai esant, kiek ji reikšminga ir iš kurių nuostatų ji kyla. \\
Pareigų registras (\en{register of obligations}) & Visos visų teisės aktų pareigos su jų
  peržiūros būsena. \\
Pasekmė (\en{consequence}) & Kiek gali kainuoti išvados nuostatos pažeidimas; kodas ją
  nustato pagal teisės akto sankcijų pakopas: tai teisės akte numatyta maksimali bauda, o ne
  prognozė. Ji rodoma tik spragai. \\
Patvirtinimas (\en{sign-off}) & Baigtos ataskaitos patvirtinimas, kurį atlieka žmogus. Jis
  atnaujina veiksmų registrą. \\
Programos dokumentas (\en{work product}) & Sistemos parengtas atitikties programos dokumentas:
  politika, valdymo struktūra su procesu ir atsakomybių matrica, rizikos registras. \\
Redakcija (\en{version}) & Vienas saugomas teisės akto teksto variantas. \\
Reikšmingumas (\en{severity}) & Kiek pareiga svarbi: kritinė, aukšta, vidutinė ar žema. Tai
  pareigos, o ne išvados savybė. \\
Sankcijų pakopa (\en{tier, sanction tier}) & Viena teisės akto sankcijų nuostatų taisyklė:
  kokie pažeidimai, kokia maksimali bauda, kokios kitos priemonės. \\
Taikytinumas (\en{applicability}) & Sprendimas, ar teisės aktas įmonei privalomas: taikoma,
  netaikoma arba neaišku (konsolė rašo \enquote{tikslintinas}). \\
Teisės aktas (\en{act}) & Teisės aktų bazėje saugomas teisinis tekstas: reglamentas, direktyva, įstatymas,
  nutarimas, įsakymas, gairės, licencijavimo puslapis. Standartas saugomas taip pat ir
  vadinamas standartu. \\
Teisės akto vertinimas (\en{assessment of an act}) & Ką ataskaita daro su teisės aktu:
  \emph{vertinamas}; \emph{dar netaikomas} --- vertinamas pasirengimo tikslu, nes privalomas
  nuo vėlesnės datos; \emph{tikslintinas} --- nevertinamas, kol įmonė neatsako į atvirus
  klausimus; \emph{taikomas netiesiogiai} --- direktyva, kuri įmonę pasiekia per ją
  perkeliantį nacionalinį įstatymą, arba teisės aktas, kuris tik nustato institucijos
  sąrangą ir įmonę pasiekia per tos institucijos sprendimus. \\
Teisės aktų bazė (\en{corpus, legal corpus}) & Visi teisės aktai su jų tekstu, saugomu
  redakcijomis ir suskaidytu į nuostatas. \\
Užduočių vykdytojas (\en{worker}) & Procesas, vykdantis užduotis. \\
Užduotis, etapas (\en{job, step}) & Foninio darbo vienetas ir jo dalys. Ataskaitą parengia
  užduotis \code{report.analyze} aštuoniais etapais, o kai vertinami ir įmonės vidaus
  reikalavimai --- devyniais; ataskaita, pradėta su apimties peržiūra, vykdoma kaip dvi tokios
  užduotys: viena prieš patvirtinimą, kita po jo. \\
Veiksmų planas (\en{action plan}) & Ataskaitos siūlomų veiksmų sąrašas; kiekviena atvira išvada
  priskirta lygiai vienam iš jų. \\
Veiksmų registras (\en{action register}) & Įmonės veiksmų sąrašas, išliekantis nuo vienos
  ataskaitos iki kitos. Kodas jį keičia tik tada, kai žmogus patvirtina ataskaitą arba patvirtinimas
  atšaukiamas; tarp šių momentų kiekvieną veiksmą žmogus prižiūri rankiniu būdu. \\
Vektorius (\en{embedding}) & Teksto ištrauką atspindintis vektorius, leidžiantis rasti panašios
  reikšmės ištraukas. \\
Vertintojas (\en{judge}) & Antras modelio iškvietimas, kuris pirmojo išvestį tikrina pagal
  šaltinį: atskiras iškvietimas su savo instrukcija, nebūtinai kitas modelis (veikiančioje
  sistemoje visi vertintojai veikia su \model{gpt-6.1-sol}). Jis abejoja arba pritaria; nieko
  netvirtina. \\
Vidaus reikalavimas (\en{own requirement}) & Reikalavimas, kurį įmonė sau nustatė vienoje iš savo
  politikų, išskirtas tam, kad ataskaita galėtų patikrinti, ar procedūros jį įgyvendina. \\
\end{xltabular}

\section{Būsenos ir jų užrašai}\label{sec:glossary:labels}

Kodas būseną saugo angliškuoju raktu; konsolė rodo užrašą sąsajos kalba. Šis dokumentas
būsenas vadina pagal saugomus pavadinimus, o \code{unknown} vadina \enquote{neįrodyta}.

\begin{xltabular}{\textwidth}{@{}L{26mm}L{26mm}L{40mm}Y@{}}
\toprule
\hd{Objektas} & \hd{Saugoma kaip} & \hd{Užrašas lietuviškai} & \hd{Užrašas angliškai} \\
\midrule
\endfirsthead
\toprule
\hd{Objektas} & \hd{Saugoma kaip} & \hd{Užrašas lietuviškai} & \hd{Užrašas angliškai} \\
\midrule
\endhead
\bottomrule
\endlastfoot
Išvada & \code{met} & Atitinka & \en{Met} \\
 & \code{partial} & Iš dalies & \en{Partial} \\
 & \code{missing} & Neatitinka & \en{Not met} \\
 & \code{unknown} & Neįrodyta & \en{Not evidenced} \\
 & \code{not_applicable} & Netaikoma & \en{Not applicable} \\
\midrule
Reikšmingumas & \code{critical} & Kritinė & \en{Critical} \\
 & \code{high} & Aukšta & \en{High} \\
 & \code{medium} & Vidutinė & \en{Medium} \\
 & \code{low} & Žema & \en{Low} \\
\midrule
Taikytinumas & \code{applies} & Taikomas & \en{Applies} \\
 & \code{not_applies} & Netaikomas & \en{Does not apply} \\
 & \code{uncertain} & Tikslintinas & \en{To be confirmed} \\
\midrule
Ataskaita & \code{queued} & Laukia eilėje & \en{Queued} \\
 & \code{running} & Vykdoma & \en{Running} \\
 & \code{scope_review} & Laukia apimties patvirtinimo & \en{Awaiting scope approval} \\
 & \code{done} & Baigta & \en{Finished} \\
 & \code{failed} & Nepavyko & \en{Failed} \\
 & \code{cancelled} & Atšaukta & \en{Cancelled} \\
\midrule
Pareiga & \code{proposed} & Pasiūlyta & \en{Proposed} \\
 & \code{approved} & Patvirtinta & \en{Approved} \\
 & \code{rejected} & Atmesta & \en{Rejected} \\
\midrule
Vertintojo verdiktas & \code{supported} & Pagrįsta & \en{Supported} \\
 & \code{partly} & Iš dalies & \en{Partly} \\
 & \code{unsupported} & Nepagrįsta & \en{Unsupported} \\
\midrule
Veiksmas & \code{open} & Atviras & \en{Open} \\
 & \code{in_progress} & Vykdomas & \en{In progress} \\
 & (išvestinė) & Vėluojantis & \en{Late} \\
 & \code{done} & Įvykdytas, laukia patikros & \en{Done, awaiting re-test} \\
 & \code{closed} & Uždarytas & \en{Closed} \\
\midrule
Programos dokumentas & \code{queued} & Laukia eilėje & \en{Queued} \\
 & \code{running} & Generuojama & \en{Generating} \\
 & \code{draft} & Projektas & \en{Draft} \\
 & \code{approved} & Patvirtinta & \en{Approved} \\
 & \code{failed} & Nepavyko & \en{Failed} \\
 & \code{cancelled} & Atšaukta & \en{Cancelled} \\
\end{xltabular}

\section{Teisės aktų santrumpos}\label{sec:glossary:acts}

Sistema teisės aktą vadina trumpuoju pavadinimu; Lietuvos įstatymo trumpasis pavadinimas yra
santrumpa, kurią vartoja Lietuvos teisininkai. Visus teisės aktus su jų reguliavimo dalyku
išvardija \cref{app:acts}.

\begin{xltabular}{\textwidth}{@{}L{20mm}YY@{}}
\toprule
\hd{Trumpinys} & \hd{Lietuviškai} & \hd{Angliškai} \\
\midrule
\endfirsthead
\toprule
\hd{Trumpinys} & \hd{Lietuviškai} & \hd{Angliškai} \\
\midrule
\endhead
\bottomrule
\endlastfoot
ADTAĮ & Asmens duomenų teisinės apsaugos įstatymas & \en{Law on the Legal Protection of Personal Data} \\
BĮ & Bankų įstatymas & \en{Law on Banks} \\
EPEPĮĮ & Elektroninių pinigų ir elektroninių pinigų įstaigų įstatymas & \en{Law on Electronic Money and Electronic Money Institutions} \\
FĮĮ & Finansų įstaigų įstatymas & \en{Law on Financial Institutions} \\
KRĮ & Kriptoturto rinkų įstatymas & \en{Law on Markets in Crypto-Assets} \\
KSĮ & Kibernetinio saugumo įstatymas & \en{Law on Cyber Security} \\
KUĮ & Kredito unijų įstatymas & \en{Law on Credit Unions} \\
LBĮ & Lietuvos banko įstatymas & \en{Law on the Bank of Lithuania} \\
MĮ & Mokėjimų įstatymas & \en{Law on Payments} \\
MĮĮ & Mokėjimo įstaigų įstatymas & \en{Law on Payment Institutions} \\
NTSKĮ & Su nekilnojamuoju turtu susijusio kredito įstatymas & \en{Law on Credit Related to Immovable Property} \\
PAĮ & Pranešėjų apsaugos įstatymas & \en{Law on the Protection of Whistleblowers} \\
PPTFPĮ & Pinigų plovimo ir teroristų finansavimo prevencijos įstatymas & \en{Law on the Prevention of Money Laundering and Terrorist Financing} \\
VKĮ & Vartojimo kredito įstatymas & \en{Law on Consumer Credit} \\
\midrule
AMLAR & Reglamentas (ES) 2024/1620, kuriuo įsteigiama Kovos su pinigų plovimu ir teroristų finansavimu institucija & \en{Regulation (EU) 2024/1620 establishing the Authority for Anti-Money Laundering and Countering the Financing of Terrorism} \\
AMLD6 & Direktyva (ES) 2024/1640 dėl mechanizmų, kuriuos valstybės narės turi įdiegti kovai su pinigų plovimu ir teroristų finansavimu & \en{Directive (EU) 2024/1640 on the mechanisms Member States put in place against money laundering and terrorist financing} \\
AMLR & Reglamentas (ES) 2024/1624 dėl pinigų plovimo ir teroristų finansavimo prevencijos & \en{Regulation (EU) 2024/1624 on the prevention of money laundering and terrorist financing} \\
CCD2 & Direktyva (ES) 2023/2225 dėl vartojimo kredito sutarčių & \en{Directive (EU) 2023/2225 on credit agreements for consumers} \\
CRA & Reglamentas (ES) 2024/2847, Kibernetinio atsparumo aktas & \en{Regulation (EU) 2024/2847, the Cyber Resilience Act} \\
DORA & Reglamentas (ES) 2022/2554 dėl finansų sektoriaus skaitmeninės veiklos atsparumo & \en{Regulation (EU) 2022/2554 on digital operational resilience for the financial sector} \\
DSA & Reglamentas (ES) 2022/2065, Skaitmeninių paslaugų aktas & \en{Regulation (EU) 2022/2065, the Digital Services Act} \\
ECSPR & Reglamentas (ES) 2020/1503 dėl Europos sutelktinio finansavimo paslaugų teikėjų & \en{Regulation (EU) 2020/1503 on European crowdfunding service providers} \\
eIDAS~2 & Reglamentas (ES) Nr.~910/2014 dėl elektroninės atpažinties ir elektroninių operacijų patikimumo užtikrinimo paslaugų, iš dalies pakeistas 2024~m. & \en{Regulation (EU) No 910/2014 on electronic identification and trust services, as amended in 2024} \\
EMD2 & Direktyva 2009/110/EB dėl elektroninių pinigų įstaigų & \en{Directive 2009/110/EC on electronic money institutions} \\
GDPR & Reglamentas (ES) 2016/679, Bendrasis duomenų apsaugos reglamentas & \en{Regulation (EU) 2016/679, the General Data Protection Regulation} \\
MAR & Reglamentas (ES) Nr.~596/2014 dėl piktnaudžiavimo rinka & \en{Regulation (EU) No 596/2014 on market abuse} \\
MiCA & Reglamentas (ES) 2023/1114 dėl kriptoturto rinkų & \en{Regulation (EU) 2023/1114 on markets in crypto-assets} \\
NIS2 & Direktyva (ES) 2022/2555 dėl aukšto bendro kibernetinio saugumo lygio visoje Sąjungoje & \en{Directive (EU) 2022/2555 on a high common level of cybersecurity across the Union} \\
PSD2 & Direktyva (ES) 2015/2366 dėl mokėjimo paslaugų vidaus rinkoje & \en{Directive (EU) 2015/2366 on payment services in the internal market} \\
TFR & Reglamentas (ES) 2023/1113 dėl informacijos, teikiamos kartu su lėšų ir tam tikro kriptoturto pervedimais & \en{Regulation (EU) 2023/1113 on information accompanying transfers of funds and certain crypto-assets} \\
\end{xltabular}

\section{Kitos santrumpos}\label{sec:glossary:other}

\begin{xltabular}{\textwidth}{@{}L{24mm}Y@{}}
\toprule
\hd{Santrumpa} & \hd{Reikšmė} \\
\midrule
\endfirsthead
\toprule
\hd{Santrumpa} & \hd{Reikšmė} \\
\midrule
\endhead
\bottomrule
\endlastfoot
CASP & Kriptoturto paslaugų teikėjas (angl. \en{crypto-asset service provider}). \\
CELEX & Europos Sąjungos teisės dokumento identifikatorius (pavyzdžiui, \code{32023R1114} ---
  MiCA). \\
CELLAR & Europos Sąjungos leidinių biuro saugykla, iš kurios sistema atsisiunčia ES teisės
  aktus. \\
EBA & Europos bankininkystės institucija (angl. \en{European Banking Authority}). \\
e-TAR & Teisės aktų registras --- Lietuvos teisės aktų duomenų bazė, iš kurios sistema atsisiunčia
  Lietuvos įstatymus, Lietuvos banko nutarimus ir FNTT įsakymą; Lietuvos banko gairės ir
  licencijavimo puslapiai paimami iš paties banko svetainės. \\
FNTT & Finansinių nusikaltimų tyrimo tarnyba. \\
HNSW & Hierarchinis naršomas mažojo pasaulio grafas (angl. \en{hierarchical navigable small
  world}): indekso rūšis, naudojama artimiausių kaimynų paieškai tarp vektorių. \\
ISMS & Informacijos saugumo valdymo sistema (angl. \en{information security management
  system}). \\
ITS, RTS & Techniniai įgyvendinimo ir techniniai reguliavimo standartai: Komisijos reglamentai,
  nustatantys, kaip įgyvendinama pagrindinio reglamento nuostata (\enquote{antrasis lygmuo}). \\
LB & Lietuvos bankas, finansų rinkos priežiūros institucija. \\
MCP & Modelio konteksto protokolas (angl. \en{Model Context Protocol}): protokolas, per kurį
  pokalbių asistentas gali naudoti įrankius, kuriuos siūlo administratoriaus užregistruotas
  išorinis serveris. \\
ML/TF & Pinigų plovimas ir teroristų finansavimas (angl. \en{money laundering and terrorist
  financing}). \\
RACI & Atsakingas, atskaitingas, konsultuojamas, informuojamas (angl. \en{responsible,
  accountable, consulted, informed}): keturi būdai, kuriais rolė dalyvauja proceso etape. \\
RRF & Atvirkštinių rangų sujungimas (angl. \en{reciprocal rank fusion}): būdas, kuriuo vektorinės
  paieškos ir viso teksto paieškos rezultatai sujungiami į vieną reitingą. \\
SNI & Serverio vardo nurodymas (angl. \en{Server Name Indication}): serverio
  vardas, kurį klientas nurodo užmegzdamas šifruotą ryšį. Vienintelis sistemos kelias į
  internetą ryšius leidžia pagal šį vardą. \\
\end{xltabular}
